\documentclass[tikz,border=10pt]{standalone}
\usepackage{ctex}
\usepackage{tikz}
\usetikzlibrary{arrows.meta,positioning,fit,backgrounds}

\begin{document}
\begin{tikzpicture}[
  font=\small,
  >=Latex,
  blk/.style={draw=#1!70!black, fill=#1!10, rounded corners=2.2mm, very thick, minimum width=28mm, minimum height=10mm, align=center},
  sig/.style={->, very thick, draw=black!75},
  box/.style={draw=black!25, rounded corners=3mm, line width=0.8pt}
]

\node[blk=blue] (raw) at (0,0) {CarHack 原始数据\\CSV/TXT};
\node[blk=teal] (prep) at (3.6,0) {预处理\\64x9 滑窗构造};
\node[blk=orange] (train) at (7.2,0) {训练\\CAN\_CNN};
\node[blk=purple] (eval) at (10.8,0) {评估与选择\\best\_model.pth};
\node[blk=green!70!black] (deploy) at (14.4,0) {在线部署\\ml\_bridge 推理};

\draw[sig] (raw) -- (prep);
\draw[sig] (prep) -- (train);
\draw[sig] (train) -- (eval);
\draw[sig] (eval) -- (deploy);

\node[draw=black!35, rounded corners=2mm, fill=yellow!12, align=left, inner sep=4pt] at (7.2,-1.8) {
关键约束：部署端必须复用训练期 `dt\_max\_ms`，\\
保证第 9 列 \(\Delta t\) 归一化分布一致。
};

\begin{scope}[on background layer]
  \node[box, fit=(raw)(prep)(train)(eval)(deploy), inner sep=7mm, label={[yshift=2mm]\large\bfseries CAN-CNN(64x9) 端到端流程}] {};
\end{scope}

\end{tikzpicture}
\end{document}
